\chapter*{\centerline{Abstract}}
\markboth{\MakeUppercase{Abstract}}{}
\iftoggle{fulltoc}{
  \addcontentsline{toc}{chapter}{Abstract}
}{}

Biomimetic robots, robots whose biomechanics are modeled on human or animal anatomy, offer a platform for studying neuromechanical control through experiments that would be impractical or unethical with living subjects. This dissertation advances a long-running effort at Portland State University to build a humanoid robot, from the trunk down, actuated by pneumatic artificial muscles and controlled by a synthetic neural network. The work addresses the foundational question of that effort: can pneumatic artificial muscle actuators be characterized and arranged so that a biomimetic robot reproduces human joint force capabilities?

First, the isometric force of Festo braided pneumatic actuators (BPAs) with uninflated diameters of \qtylist{10;20}{\mm} was characterized over a wide range of resting lengths, contractions, and pressures using a custom test jig. The analysis produced novel equations for the maximum isometric force at \qty{620}{\kPa} as a function of resting length --- a previously unreported dependency that the manufacturer's tool predicts with the wrong trend --- and a three-coefficient nondimensional force surface that predicts isometric force from relative strain and relative pressure with high fidelity (adjusted $R^2 > 0.99$), correcting the large force over-predictions of existing tools and models at short resting lengths.

Second, the actuator model was carried to the joint level through isometric knee torque experiments on two test stands: a simplified pinned knee and a biomimetic four-bar knee with a migrating instantaneous center of rotation. A hybrid torque calculation method isolated the sources of prediction error, revealing constant length offsets, series compliance in brackets and artificial tendons, and a loss of usable actuator length when BPAs wrap around the joint. Correction terms for each effect were identified through multiobjective optimization and substantially improved torque prediction across flexor and extensor configurations. The biomimetic knee driven by optimized BPA placements was shown to meet or exceed human isometric knee torque over most of the range of motion.

Third, the dissertation defines the path from this validated biomechanical plant to neuromechanical control: conversion of the lab's modified OpenSim leg model to MuJoCo via MyoConverter, a synthetic nervous system implemented in SNS-Toolbox built around a two-layer central pattern generator with per-leg rhythm generation and pattern formation and left--right coordination, proprioceptive (Ia, Ib, II) and heel/toe contact feedback, and verification by tonic-stimulation tests against the locomotion literature. 
